\documentclass[11pt,a4paper]{article}
% Common preamble for the hint sheets (\input from each file)
\usepackage[utf8]{inputenc}
\usepackage[T1]{fontenc}
\usepackage{lmodern}
\usepackage{amsmath,amssymb,amsthm,mathtools}
\usepackage[a4paper,margin=2.3cm]{geometry}
\usepackage{enumitem}
\usepackage{xcolor}
\usepackage[hidelinks]{hyperref}
\renewcommand{\leq}{\leqslant}
\renewcommand{\le}{\leqslant}
\renewcommand{\geq}{\geqslant}
\renewcommand{\ge}{\geqslant}
\setlength{\parindent}{0pt}
\setlength{\parskip}{0.35em}
% \hintstitle{Sheet title}{Class work / Homework}
\newcommand{\hintstitle}[2]{\begin{center}{\LARGE\bfseries Hints: #1\par}\vspace{0.5em}{\large #2\par}\end{center}
\noindent\textit{The hints to each problem are cumulative. Read one, try again, and open the next only if you are still stuck.}\par\medskip\hrule\medskip}
\newcommand{\problem}[1]{\par\bigskip{\large\bfseries Problem #1}\par\nopagebreak\smallskip}
\newcommand{\subpart}[1]{\par\smallskip\textbf{(#1)}\par\nopagebreak}
\newcommand{\hint}[1]{\par\smallskip\noindent\textcolor{black!70}{\textbf{Hint #1.}}\enspace}
\newcommand{\note}[1]{\par\smallskip\noindent\textit{Note.} #1\par}
\providecommand{\lt}{<}
\providecommand{\gt}{>}
\providecommand{\tightlist}{\setlength{\itemsep}{0pt}\setlength{\parskip}{0pt}}
\emergencystretch=3em

\begin{document}
\hintstitle{Limits -- 3. Asymptotics of sequences}{Class work}

\problem{1}
\hint{1} The sequence consists of blocks $1,2,\dots,m$. Compute the sum $S_n=\sum_{k\le n}a_k$ exactly when $n=T_m=\frac{m(m+1)}2$ is the end of the $m$-th block.
\hint{2} $S_{T_m}=\sum_{j=1}^m\frac{j(j+1)}2=\frac{m(m+1)(m+2)}6$. For $T_m\le n<T_{m+1}$ one has $S_{T_m}\le S_n\le S_{T_{m+1}}$, and $m\sim\sqrt{2n}$.
\hint{3} So $S_n\sim\frac{m^3}6\sim\frac{(2n)^{3/2}}6=\frac{\sqrt2}3\,n^{3/2}$. The limit is finite and positive only for $\alpha=\frac32$; answer: $(\alpha,\beta)=\bigl(\frac32,\frac{\sqrt2}3\bigr)$.

\problem{2}
\note{The condition $0<x_0<x_1$ is impossible, since $x_1=x_0-x_0^2<x_0$. The intended condition is $0<x_0<1$; the hints use it.}
\hint{1} Show that $0<x_n<1$ and $x_n$ decreases, so $x_n\to L$ with $L=L-L^2$, i.e.\ $L=0$.
\hint{2} Telescoping (Fact 1 with $p=2$): $\dfrac1{x_{n+1}}-\dfrac1{x_n}=\dfrac1{1-x_n}\to1$. By Stolz--Ces\`aro (or Ces\`aro means) $\frac1{nx_n}\to1$, so $\lim nx_n=1$.

\problem{3}
\hint{1} First $x_1\in(0,1]$ and $x_n$ decreases to $0$. Expand $\dfrac1{\sin^2x}=\dfrac1{x^2}+\dfrac13+\dfrac{x^2}{15}+O(x^4)$, so $\frac1{x_{n+1}^2}-\frac1{x_n^2}\to\frac13$.
\hint{2} By Stolz, $\frac1{nx_n^2}\to\frac13$, i.e.\ $\sqrt n\,x_n\to\sqrt3$.
\hint{3} For the next term, sum the expansion: $\frac1{x_n^2}=\frac n3+\frac1{15}\sum_{k<n}x_k^2+O(1)$, and $x_k^2=\frac3k+o(\frac1k)$ gives $\frac1{x_n^2}=\frac n3+\frac{\ln n}5+o(\ln n)$ (in fact $+\,C+o(1)$, since the error terms are summable).
\hint{4} Therefore $x_n=\sqrt{\dfrac3n}\Bigl(1-\dfrac{3\ln n}{10\,n}+O\bigl(\tfrac1n\bigr)\Bigr)=\sqrt{\dfrac3n}-\dfrac{3\sqrt3}{10}\,\dfrac{\ln n}{n^{3/2}}+O\bigl(n^{-3/2}\bigr)$.

\problem{4}
\note{As printed the answer is $\lim x_n^3\ln n=0$; the natural quantity with a finite nonzero limit is $x_n^2\ln n\to3$. The hints derive both.}
\hint{1} Put $y_n=x_{n+1}^3=\frac1n\sum_{k\le n}\sin^3x_k$. Then $(n+1)y_{n+1}=ny_n+\sin^3x_{n+1}$, i.e.\ $y_{n+1}-y_n=\dfrac{\sin^3x_{n+1}-x_{n+1}^3}{n+1}<0$. Deduce that $y_n\downarrow0$ (if $y_n\to L>0$, the decrements are $\ge c/n$ and sum to $\infty$).
\hint{2} Since $\sin^3t-t^3=-\frac{t^5}2+O(t^7)$, we get $y_{n+1}=y_n-\dfrac{y_n^{5/3}}{2(n+1)}(1+o(1))$. This is a neutral fixed point with an extra factor $\frac1n$: telescope in the power $y^{-2/3}$.
\hint{3} $y_{n+1}^{-2/3}-y_n^{-2/3}=\dfrac{1+o(1)}{3(n+1)}$, so by Stolz with $b_n=\ln n$: $y_n^{-2/3}\sim\frac{\ln n}3$. Hence $x_n^2\sim\frac3{\ln n}$, so $x_n^2\ln n\to3$ and $x_n^3\ln n\sim 3^{3/2}(\ln n)^{-1/2}\to0$.

\problem{5}
\hint{1} No continuity is assumed, so use order instead of the intermediate value theorem: consider $A=\{x\in[a,b]: f(x)\ge x\}$, which contains $a$, and $\xi=\sup A$.
\hint{2} For $x\in A$: $x\le f(x)\le f(\xi)$, so $f(\xi)$ is an upper bound of $A$ and $\xi\le f(\xi)$. Then $f(f(\xi))\ge f(\xi)$, so $f(\xi)\in A$ and $f(\xi)\le\xi$.
\hint{3} Example: $f(x)=1-x$ on $[0,1]$ has the unique fixed point $\frac12$, but $x_0=0$ gives $0,1,0,1,\dots$
\end{document}
